% Options for packages loaded elsewhere
\PassOptionsToPackage{unicode}{hyperref}
\PassOptionsToPackage{hyphens}{url}
\documentclass[aspectratio=169,ignorenonframetext]{beamer}
\usepackage{ma2003b-beamer}
\hypersetup{
  pdftitle={Factor Analysis \& PCA},
  pdfauthor={Juliho Castillo Colmenares},
  hidelinks,
  pdfcreator={LaTeX via pandoc}}

\title{\texorpdfstring{Factor Analysis \& PCA}{Factor Analysis \& PCA}}
\author{\texorpdfstring{Juliho Castillo
Colmenares}{Juliho Castillo Colmenares}}
\subtitle{Latent Variable Modeling and Dimensionality Reduction}
\institute{Tec de Monterrey}
\date{}

\begin{document}
\frame{\titlepage}

\begin{frame}{Today's Agenda}
\begin{enumerate}
\item
  Introduction to Dimensionality Reduction
\item
  Principal Component Analysis (PCA) Review
\item
  The Common Factor Model (FA)
\item
  Factor Extraction \& Communalities
\item
  Factor Rotation (Orthogonal vs. Oblique)
\item
  PCA vs. Factor Analysis: Methodological Comparison
\item
  Suitability Testing (KMO \& Bartlett's)
\item
  Practical Python Implementation \& Educational Case Study
\end{enumerate}
\end{frame}

\begin{frame}{Motivation: The Curse of Dimensionality}
\textbf{Why reduce dimensions in multivariate data?}

\begin{itemize}
\item
  \textbf{High Collinearity:} Redundant predictors destabilize
  regression and inference.
\item
  \textbf{Interpretability:} Condense dozens of observed indicators into
  underlying constructs.
\item
  \textbf{Visualization:} Project \(p\)-dimensional data onto 2D or 3D
  coordinate planes.
\item
  \textbf{Noise Reduction:} Separate shared systematic signal from
  item-specific measurement noise.
\end{itemize}

{\def\LTcaptype{none} % do not increment counter
\begin{longtable}[]{@{}
  >{\raggedright\arraybackslash}p{(\linewidth - 2\tabcolsep) * \real{0.5000}}
  >{\raggedright\arraybackslash}p{(\linewidth - 2\tabcolsep) * \real{0.5000}}@{}}
\toprule\noalign{}
\endhead
\begin{minipage}[t]{\linewidth}\raggedright
\textbf{Principal Component Analysis (PCA)}

\begin{itemize}
\item
  Data-driven / Variance maximization
\item
  Linear combinations of all variance
\item
  Orthogonal components by design
\end{itemize}
\end{minipage} & \begin{minipage}[t]{\linewidth}\raggedright
\textbf{Exploratory Factor Analysis (EFA)}

\begin{itemize}
\item
  Theory-driven / Latent construct modeling
\item
  Partitions \textbf{common} vs. \textbf{unique} variance
\item
  Rotatable for domain interpretability
\end{itemize}
\end{minipage} \\
\bottomrule\noalign{}
\end{longtable}
}
\end{frame}

\begin{frame}{Principal Component Analysis (PCA)}
\textbf{Objective:} Transform \(p\) correlated variables
\(X_{1},\ldots,X_{p}\) into uncorrelated components
\(Y_{1},\ldots,Y_{p}\) ordered by variance.

\textbf{Mathematical Foundation:} Given covariance matrix \(\mathbf{S}\)
or correlation matrix \(\mathbf{R}\):
\[\mathbf{S} = \mathbf{V}\mathbf{\Lambda}\mathbf{V}^{T}\]

where
\(\mathbf{\Lambda} = \operatorname{diag}(\lambda_{1},\lambda_{2},\ldots,\lambda_{p})\)
with
\(\lambda_{1} \geq \lambda_{2} \geq \ldots \geq \lambda_{p} \geq 0\),
and
\(\mathbf{V} = \left\lbrack \mathbf{v}_{1},\ldots,\mathbf{v}_{p} \right\rbrack\)
are orthonormal eigenvectors.

\textbf{Principal Components:} \[\begin{array}{r}
Y_{i} = \mathbf{v}_{i}^{T}\left( \mathbf{X} - \overline{\mathbf{x}} \right),\quad\operatorname{Var}(Y_{i}) = \lambda_{i},\quad\operatorname{Cov}(Y_{i},Y_{j}) = 0 \\
(i \neq j)
\end{array}\]
\end{frame}

\begin{frame}{Component Retention Criteria}
{\def\LTcaptype{none} % do not increment counter
\begin{longtable}[]{@{}
  >{\raggedright\arraybackslash}p{(\linewidth - 2\tabcolsep) * \real{0.5000}}
  >{\raggedright\arraybackslash}p{(\linewidth - 2\tabcolsep) * \real{0.5000}}@{}}
\toprule\noalign{}
\endhead
\begin{minipage}[t]{\linewidth}\raggedright
\textbf{1. Kaiser-Guttman Criterion:}

\begin{itemize}
\item
  Retain components with eigenvalues \(\lambda_{i} > 1\) (for
  standardized \(\mathbf{R}\)).
\item
  Explains more variance than a single original variable.
\end{itemize}

\textbf{2. Cumulative Variance Threshold:}

\begin{itemize}
\item
  Retain enough components to explain \(70\% - 85\%\) of total variance:
\end{itemize}

\[\frac{\sum_{i = 1}^{k}\lambda_{i}}{\sum_{j = 1}^{p}\lambda_{j}} \geq \text{ threshold }\]
\end{minipage} & \begin{minipage}[t]{\linewidth}\raggedright
\textbf{3. Scree Plot (Cattell):}

\begin{itemize}
\item
  Plot \(\lambda_{i}\) vs. component index \(i\).
\item
  Identify the ``elbow'' where the curve levels off.
\end{itemize}

\textbf{4. Parallel Analysis:}

\begin{itemize}
\item
  Compare empirical eigenvalues against eigenvalues simulated from
  random data of identical dimensions.
\end{itemize}
\end{minipage} \\
\bottomrule\noalign{}
\end{longtable}
}
\end{frame}

\begin{frame}{The Common Factor Model}
\textbf{Hypothesis:} Observed variables \(X_{1},\ldots,X_{p}\) are
generated by \(k < p\) unobserved \textbf{latent factors}
\(F_{1},\ldots,F_{k}\) plus unique errors.

\textbf{Vector Formulation:}
\[\mathbf{X} - \mathbf{\mu} = \mathbf{\Lambda}\mathbf{F} + \mathbf{\epsilon}\]

where:

\begin{itemize}
\item
  \(\mathbf{X} = \left\lbrack X_{1},\ldots,X_{p} \right\rbrack^{T}\):
  Observed variables vector.
\item
  \(\mathbf{\Lambda} \in \mathbb{R}^{p \times k}\): Factor loadings
  matrix
  (\(\lambda_{ji} = \text{ loading of  }X_{j}\text{  on  }F_{i}\)).
\item
  \(\mathbf{F} = \left\lbrack F_{1},\ldots,F_{k} \right\rbrack^{T}\):
  Latent common factor scores
  (\(\mathbb{E}\left\lbrack \mathbf{F} \right\rbrack = \mathbf{0},\operatorname{Cov}(\mathbf{F}) = \mathbf{I}_{k}\)).
\item
  \(\mathbf{\epsilon} = \left\lbrack \epsilon_{1},\ldots,\epsilon_{p} \right\rbrack^{T}\):
  Specific / unique error vector
  (\(\operatorname{Cov}(\mathbf{\epsilon}) = \mathbf{\Psi} = \operatorname{diag}(\psi_{1}^{2},\ldots,\psi_{p}^{2})\)).
\end{itemize}
\end{frame}

\begin{frame}{Variance Partitioning: Communality \& Uniqueness}
For standardized variables, total variance of \(X_{j}\) equals 1:
\[\operatorname{Var}(X_{j}) = \underset{h_{j}^{2}\text{  (Communality)}}{\underbrace{\sum_{i = 1}^{k}\lambda_{ji}^{2}}} + \underset{\text{ Uniqueness}}{\underbrace{\psi_{j}^{2}}} = 1\]

\textbf{Key Metrics:}

\begin{itemize}
\item
  \textbf{Communality (\(h_{j}^{2}\)):} Proportion of \(X_{j}\)`s
  variance explained by the shared \(k\) common factors
  (\(h_{j}^{2} \in \lbrack 0,1\rbrack\)).
\item
  \textbf{Uniqueness (\(\psi_{j}^{2} = 1 - h_{j}^{2}\)):} Variance
  specific to variable \(X_{j}\) (specific factor + measurement error).
\item
  \textbf{Reproduced Correlation Matrix:}
\end{itemize}

\[\hat{\mathbf{R}} = \mathbf{\Lambda}\mathbf{\Lambda}^{T} + \mathbf{\Psi}\]
\end{frame}

\begin{frame}{Factor Extraction Methods}
\textbf{1. Principal Axis Factoring (PAF):}

\begin{itemize}
\item
  Iteratively estimates communalities along the diagonal of
  \(\mathbf{R} - \hat{\mathbf{\Psi}}\).
\item
  Non-parametric; does not require multivariate normality.
\end{itemize}

\textbf{2. Maximum Likelihood (ML) Extraction:}

\begin{itemize}
\item
  Assumes multivariate normal data:
  \(\mathbf{X} \sim \mathcal{N}_{p\left( \mathbf{\mu},\mathbf{\Sigma} \right)}\).
\item
  Maximizes the log-likelihood of observing sample covariance
  \(\mathbf{S}\).
\item
  Provides statistical hypothesis tests (\(\chi^{2}\) goodness-of-fit)
  for the number of factors \(k\).
\end{itemize}
\end{frame}

\begin{frame}{Factor Rotation: Achieving Simple Structure}
\textbf{The Rotation Problem:} Initial factor extraction is
mathematically indeterminate up to an orthogonal matrix transformation
\(\mathbf{T}\)
(\(\mathbf{\Lambda}^{\ast} = \mathbf{\Lambda}\mathbf{T}\)).

\textbf{Thurstone's Simple Structure:} Each variable should have large
loadings on \textbf{only one} factor and near-zero loadings on others.

{\def\LTcaptype{none} % do not increment counter
\begin{longtable}[]{@{}
  >{\raggedright\arraybackslash}p{(\linewidth - 6\tabcolsep) * \real{0.2000}}
  >{\raggedright\arraybackslash}p{(\linewidth - 6\tabcolsep) * \real{0.2667}}
  >{\raggedright\arraybackslash}p{(\linewidth - 6\tabcolsep) * \real{0.2667}}
  >{\raggedright\arraybackslash}p{(\linewidth - 6\tabcolsep) * \real{0.2667}}@{}}
\toprule\noalign{}
\begin{minipage}[b]{\linewidth}\raggedright
\textbf{Method}
\end{minipage} & \begin{minipage}[b]{\linewidth}\raggedright
\textbf{Type}
\end{minipage} & \begin{minipage}[b]{\linewidth}\raggedright
\textbf{Factor Correlation}
\end{minipage} & \begin{minipage}[b]{\linewidth}\raggedright
\textbf{Best For}
\end{minipage} \\
\midrule\noalign{}
\endhead
Varimax & Orthogonal & \(\operatorname{Cov}(F_{i},F_{j}) = 0\) &
Distinct, independent concepts \\
Quartimax & Orthogonal & \(\operatorname{Cov}(F_{i},F_{j}) = 0\) &
Simplifying rows (variables) \\
Promax & Oblique & \(\operatorname{Cov}(F_{i},F_{j}) \neq 0\) &
Realistic correlated constructs \\
Direct Oblimin & Oblique & \(\operatorname{Cov}(F_{i},F_{j}) \neq 0\) &
Tunable correlation (\(\delta\)) \\
\bottomrule\noalign{}
\end{longtable}
}
\end{frame}

\begin{frame}{Methodological Comparison: PCA vs. Factor Analysis}
{\def\LTcaptype{none} % do not increment counter
\begin{longtable}[]{@{}
  >{\raggedright\arraybackslash}p{(\linewidth - 4\tabcolsep) * \real{0.2728}}
  >{\raggedright\arraybackslash}p{(\linewidth - 4\tabcolsep) * \real{0.3637}}
  >{\raggedright\arraybackslash}p{(\linewidth - 4\tabcolsep) * \real{0.3637}}@{}}
\toprule\noalign{}
\begin{minipage}[b]{\linewidth}\raggedright
\textbf{Dimension}
\end{minipage} & \begin{minipage}[b]{\linewidth}\raggedright
\textbf{Principal Component Analysis (PCA)}
\end{minipage} & \begin{minipage}[b]{\linewidth}\raggedright
\textbf{Factor Analysis (FA)}
\end{minipage} \\
\midrule\noalign{}
\endhead
Objective & Maximize total variance explained & Model latent causal
constructs \\
Variance Analyzed & Total variance (\(1.0\) on diagonal) & Common
variance (\(h_{j}^{2}\) on diagonal) \\
Error Term & No explicit error modeling & Explicit uniqueness
\(\mathbf{\epsilon} \sim \left( 0,\mathbf{\Psi} \right)\) \\
Rotation & Uncommon (alters variance maximization) & Standard
(Varimax/Promax) \\
Assumptions & Minimal distribution assumptions & Multivariate normality
for ML \\
Best Used When & Data compression, multicollinearity fix & Construct
validation, survey scales \\
\bottomrule\noalign{}
\end{longtable}
}
\end{frame}

\begin{frame}{Pre-Analysis Suitability Checks}
Before extracting factors, evaluate correlation matrix factorability:

\textbf{1. Kaiser-Meyer-Olkin (KMO) Measure of Sampling Adequacy:}

\begin{itemize}
\item
  Quantifies proportion of variance among variables that might be
  common.
\item
  Scale: \(> 0.9\) (Marvelous), \(> 0.8\) (Meritorious), \(> 0.7\)
  (Middling), \(< 0.5\) (Unacceptable).
\end{itemize}

\textbf{2. Bartlett's Test of Sphericity:}

\begin{itemize}
\item
  \(H_{0}:\mathbf{R} = \mathbf{I}_{p}\) (Variables are completely
  uncorrelated).
\item
  \(H_{1}:\mathbf{R} \neq \mathbf{I}_{p}\).
\item
  Requires \(p < 0.05\) to justify factor extraction.
\end{itemize}
\end{frame}

\begin{frame}[fragile]{Python Implementation: End-to-End Workflow}
\begin{Shaded}
\begin{Highlighting}[]
\ImportTok{from}\NormalTok{ factor\_analyzer }\ImportTok{import}\NormalTok{ FactorAnalyzer, calculate\_kmo, calculate\_bartlett\_sphericity}
\ImportTok{from}\NormalTok{ sklearn.preprocessing }\ImportTok{import}\NormalTok{ StandardScaler}
\ImportTok{import}\NormalTok{ pandas }\ImportTok{as}\NormalTok{ pd}

\CommentTok{\# 1. Standardize indicators}
\NormalTok{X\_scaled }\OperatorTok{=}\NormalTok{ StandardScaler().fit\_transform(df)}

\CommentTok{\# 2. Check suitability}
\NormalTok{kmo\_all, kmo\_model }\OperatorTok{=}\NormalTok{ calculate\_kmo(X\_scaled)}
\NormalTok{chi2\_val, p\_val }\OperatorTok{=}\NormalTok{ calculate\_bartlett\_sphericity(X\_scaled)}
\BuiltInTok{print}\NormalTok{(}\SpecialStringTok{f"KMO: }\SpecialCharTok{\{}\NormalTok{kmo\_model}\SpecialCharTok{:.3f\}}\SpecialStringTok{ | Bartlett p{-}value: }\SpecialCharTok{\{}\NormalTok{p\_val}\SpecialCharTok{:.4e\}}\SpecialStringTok{"}\NormalTok{)}

\CommentTok{\# 3. Fit Factor Analysis with Promax rotation (oblique {-}{-} the safer default}
\CommentTok{\#    when latent traits could plausibly correlate)}
\NormalTok{fa }\OperatorTok{=}\NormalTok{ FactorAnalyzer(n\_factors}\OperatorTok{=}\DecValTok{3}\NormalTok{, rotation}\OperatorTok{=}\StringTok{"promax"}\NormalTok{, method}\OperatorTok{=}\StringTok{"principal"}\NormalTok{)}
\NormalTok{fa.fit(X\_scaled)}

\CommentTok{\# 4. Inspect factor structure, including the factor correlation matrix Phi}
\NormalTok{loadings }\OperatorTok{=}\NormalTok{ pd.DataFrame(fa.loadings\_, index}\OperatorTok{=}\NormalTok{df.columns, columns}\OperatorTok{=}\NormalTok{[}\StringTok{"F1"}\NormalTok{, }\StringTok{"F2"}\NormalTok{, }\StringTok{"F3"}\NormalTok{])}
\NormalTok{communalities }\OperatorTok{=}\NormalTok{ pd.Series(fa.get\_communalities(), index}\OperatorTok{=}\NormalTok{df.columns)}
\NormalTok{phi }\OperatorTok{=}\NormalTok{ pd.DataFrame(fa.phi\_, columns}\OperatorTok{=}\NormalTok{[}\StringTok{"F1"}\NormalTok{, }\StringTok{"F2"}\NormalTok{, }\StringTok{"F3"}\NormalTok{], index}\OperatorTok{=}\NormalTok{[}\StringTok{"F1"}\NormalTok{, }\StringTok{"F2"}\NormalTok{, }\StringTok{"F3"}\NormalTok{])}
\end{Highlighting}
\end{Shaded}
\end{frame}

\begin{frame}{Case Study: Educational Assessment (Chapter 4)}
\textbf{Synthetic data:} 200 simulated students, generated from a known
three-factor model with correlated factors (0.20-0.30). Results below
demonstrate recovering that known structure -- not a real
construct-validation finding.

\textbf{Context:} 200 simulated students evaluated across 9 cognitive
and social performance measures:

{\def\LTcaptype{none} % do not increment counter
\begin{longtable}[]{@{}
  >{\raggedright\arraybackslash}p{(\linewidth - 4\tabcolsep) * \real{0.3333}}
  >{\raggedright\arraybackslash}p{(\linewidth - 4\tabcolsep) * \real{0.3333}}
  >{\raggedright\arraybackslash}p{(\linewidth - 4\tabcolsep) * \real{0.3333}}@{}}
\toprule\noalign{}
\endhead
\begin{minipage}[t]{\linewidth}\raggedright
\textbf{Quantitative Ability}

\begin{itemize}
\item
  MathScore
\item
  AlgebraScore
\item
  GeometryScore
\end{itemize}
\end{minipage} & \begin{minipage}[t]{\linewidth}\raggedright
\textbf{Verbal Ability}

\begin{itemize}
\item
  ReadingComp
\item
  Vocabulary
\item
  Writing
\end{itemize}
\end{minipage} & \begin{minipage}[t]{\linewidth}\raggedright
\textbf{Interpersonal Skills}

\begin{itemize}
\item
  Collaboration
\item
  Leadership
\item
  Communication
\end{itemize}
\end{minipage} \\
\bottomrule\noalign{}
\end{longtable}
}

\textbf{Results (recovering the known simulated structure):}

\begin{itemize}
\item
  KMO = \(0.799\) (Acceptable) and Bartlett's test \(p < 0.001\).
\item
  Scree plot \& Kaiser criterion retain 3 factors, matching the
  simulation.
\item
  Promax (primary) rotation yields clean simple structure and recovers
  factor correlations \(\Phi \approx 0.19 - 0.36\), consistent with the
  0.20-0.30 correlations used to generate the data. Varimax (comparison)
  gives similar loadings but forces \(\Phi = \mathbf{I}\),
  misrepresenting the known structure.
\end{itemize}
\end{frame}

\begin{frame}{Summary \& Best Practices}
\begin{itemize}
\item
  \textbf{Use PCA} for data compression, image processing, and
  pre-processing for machine learning pipelines.
\item
  \textbf{Use Factor Analysis} when discovering or confirming latent
  psychological, educational, or financial traits.
\item
  Always report KMO (\(> 0.60\)) and Bartlett's test before interpreting
  factor models.
\item
  Choose \textbf{Varimax} for independent factor scores and
  \textbf{Promax} when latent traits are expected to interact.
\item
  Check communalities (\(h^{2} > 0.40\)) to ensure observed variables
  are adequately represented.
\end{itemize}

{ \textbf{Next Chapter: Chapter 5 --- Discriminant Analysis (Supervised
Classification)} }
\end{frame}

\end{document}




